\clearpage\thispagestyle{empty}
\pdfbookmark[0]{Sources for the November 1804 Juno Report}{juno180411-sources}
\begin{center}\textbf{Sources for the November 1804 Juno Report}\end{center}
\noindent This AI translation follows Franz Xaver von Zach’s report as excerpted in \href{https://gdz.sub.uni-goettingen.de/id/PPN236007467}{Gauss’s \textit{Werke}, vol.~VI} (Göttingen, 1874), pp.~256--258. The original report, \textit{Fortgesetzte Nachrichten über den neuen Harding’schen Planeten Juno}, appeared as article XXXVIII of the \textit{Monatliche Correspondenz}, vol.~X, November 1804, \href{https://www.digitale-sammlungen.de/en/view/bsb10538602?page=477}{pp.~463}--\href{https://www.digitale-sammlungen.de/en/view/bsb10538602?page=485}{471}. The collected excerpt begins with the comparison star on \href{https://www.digitale-sammlungen.de/en/view/bsb10538602?page=480}{p.~466}. It quotes Carl Friedrich Gauss on his second determination of the elements and Heinrich Wilhelm Matthias Olbers on his hypothesis, including the aphelion shifts calculated by Barnaba Oriani. Gauss’s comparison table is on \href{https://www.digitale-sammlungen.de/en/view/bsb10538602?page=483}{p.~469}. The collected edition describes the ephemeris from \href{https://www.digitale-sammlungen.de/en/view/bsb10538602?page=484}{p.~470} in square brackets.
\par\smallskip\noindent Translation and typesetting produced by AI. Journal scans from the Bayerische Staatsbibliothek.
\clearpage
